\uestcachievements

% 以下内容采用学校撰写规范中的官方示例。
% 请仅保留攻读学位期间取得的、与学位论文相关的成果，并替换全部示例内容。
% 每项成果使用一条\item，并用\textbf{}加粗本人姓名。
% 作者超过5人时，可按学校规范写为“本人姓名(排名/作者总数)”。

{\small
\begin{enumerate}[
    label={\textup{[\arabic*]}},
    labelwidth=2em,
    labelsep=.5em,
    leftmargin=2.5em,
    align=left,
    topsep=0pt,
    partopsep=0pt,
    itemsep=0pt,
    parsep=0pt
]
    \item \textbf{Zhang M}, Li M. New memory method of impedance elements for marching-on-in-time solution of time-domain integral equation[J]. Electromagnetics, 2010, 30(5): 448-462.
    \item Zhao M, \textbf{Zhang M}, Zhao M, Jiang M, Wei M, Nie M. Domain decomposition method based on integral equation for solution of scattering from very thin, conducting cavity[J]. IEEE Transactions on Antennas and Propagation, 2014, 62(10): 5344-5348.
    \item \textbf{张某}, 李某某. 时间步进算法中阻抗矩阵的高效存储新方法[J]. 电波科学学报. 2010, 25(4): 624-631.
    \item \textbf{张某}, 李某某, 王某. 时域磁场积分方程时间步进算法后时稳定性研究[J]. 电子科技大学学报（已录用）.
    \item \textbf{Zhang M}. Parameters discussion in two-level plane wave time-domain algorithm[C]. IEEE International Workshop on Electromagnetics. Chengdu, 2012: 38-39.
    \item 李某某, 王某, 刘某, \textbf{张某}. XXX检测方法与装置. 国家技术发明奖二等奖, 中华人民共和国国务院, 2018-12-12.
    \item \textbf{张某}(3/10). XXX关键技术与装备. 军队科学技术进步奖三等奖, 中央军委军事技术委员会, 2020-12-25.
\end{enumerate}
}
